\chapter{انكسار الضوء}\label{ch:g10:refraction}

اغمس مصّاصة في \omterm{def:g10:orders-of-magnitude:scinot}{كأس} ماء: تبدو مكسورة عند خط الماء.
والمسابح دائمًا أعمق مما تبدو، والماس يقذف نارًا ملوّنة من حجر
عديم اللون، والجملة التي تقرؤها الآن عبرت على الأرجح محيطًا داخل
خيط زجاجي أرقّ من الشعرة. وقانون واحد يفسّر الأربعة جميعًا --- مساواة
بين جيبين، وجدها سنيل سنة 1621 بعد خمسة عشر قرنًا من تحديق الفيزيائيين
في جداول الزوايا دون أن يروا النمط.

\section{الانعكاس: تذكير}

في وسط شفاف متجانس، ينتقل الضوء في خطوط مستقيمة ---
أي أشعة \cref{ch:g10:light-spectra}. وعند الحد الفاصل بين وسطين،
يرتدّ جزء من الضوء (\emph{الانعكاس}) ويعبر جزء آخر
(\emph{\omterm{def:g10:refraction:refraction}{الانكسار}}). وكلاهما يخضع لقوانين هندسية دقيقة، تُذكر
باصطلاح واحد: تُقاس الزوايا من \emph{\omterm{def:g10:refraction:angles}{الناظم}}، لا
من السطح أبدًا.

\begin{definition}[الناظم وزاوية الورود]\label{def:g10:refraction:angles}
عند النقطة التي يلاقي فيها شعاعٌ سطحًا، يكون
\emph{الناظم}\index{ناظم} هو المستقيم العمودي على السطح في
تلك النقطة. و\emph{زاوية الورود}\index{زاوية الورود} $i_1$
هي الزاوية بين الشعاع الوارد والناظم؛ وتُقاس زاويتا
الانعكاس \omterm{def:g10:refraction:refraction}{والانكسار} من الناظم على النحو
نفسه.
\end{definition}

\begin{proposition}[قانون الانعكاس]\label{prop:g10:refraction:reflection}
يقع الشعاع المنعكس في المستوي الذي يحوي الشعاع الوارد
\omterm{def:g10:refraction:angles}{والناظم}، على الجهة الأخرى من \omterm{def:g10:refraction:angles}{الناظم}، وتساوي زاوية الانعكاس
\omterm{def:g10:refraction:angles}{زاوية الورود}:
\[
i_1' = i_1 .
\]
\end{proposition}
\admitted

\begin{example}[قراءة الزوايا قراءة صحيحة]\label{ex:g10:refraction:mirror}
تصيب حزمة ليزر مرآة ``بزاوية $35^\circ$'' --- مقيسة من
سطح المرآة، كما تُوصف الحزم في الغالب. \omterm{def:g10:refraction:angles}{فزاوية الورود}
$90^\circ - 35^\circ = 55^\circ$، ومن ثم يغادر الشعاع المنعكس هو أيضًا بزاوية
$55^\circ$ عن \omterm{def:g10:refraction:angles}{الناظم}، وتنعطف الحزمة بمقدار
$180^\circ - 2 \times 55^\circ = 70^\circ$. ونسيان اصطلاح
\omterm{def:g10:refraction:angles}{الناظم} أشيع أخطاء هذا الفصل؛ وهو يكلّف
$90^\circ$ في كل مرة.
\end{example}

\section{الانكسار وقرينة الانكسار}

أمسِك قلمًا في \omterm{def:g10:orders-of-magnitude:scinot}{كأس} نصف ممتلئة: يبدو مكسورًا عند السطح.
فالضوء الآتي من الجزء المغمور يغيّر اتجاهه وهو يغادر الماء،
فتعيد العين --- وهي تفترض أشعة مستقيمة --- بناء الرأس
في مكان ليس فيه.

\begin{definition}[الانكسار]\label{def:g10:refraction:refraction}
\emph{الانكسار}\index{انكسار} هو تغيّر اتجاه الضوء
وهو يعبر الحد الفاصل بين وسطين شفافين. والشعاع في
الوسط الثاني هو \emph{الشعاع المنكسر}\index{شعاع منكسر}،
وزاويته عن \omterm{def:g10:refraction:angles}{الناظم} هي زاوية الانكسار $i_2$.
\end{definition}

ويحدث \omterm{def:g10:refraction:refraction}{الانكسار} لأن للضوء سرعات مختلفة في الأوساط
المختلفة؛ ويتميز كل مادة بمقدار ما تبطّئ به الضوء.

\begin{definition}[قرينة الانكسار]\label{def:g10:refraction:index}
\emph{قرينة انكسار}\index{قرينة الانكسار} وسطٍ شفاف
هي النسبة
\[
n = \frac{c}{v},
\]
حيث $c = \qty{3.00e8}{m/s}$ سرعة الضوء في الخلاء و$v$ سرعته
في الوسط. وبما أن $v \leq c$، تحقق القرينة $n \geq 1$؛ وهي
بلا \omterm{def:g10:orders-of-magnitude:unit}{وحدة}. وإليك بعض القيم المفيدة:
\begin{center}
\begin{tabular}{lccccc}
الوسط & الهواء & الماء & البلكسيغلاس & الزجاج & الماس \\
$n$ & $1.0003$ & $1.33$ & $1.49$ & $1.50$ & $2.42$ \\
\end{tabular}
\end{center}
\end{definition}

\begin{example}[سرعة الضوء في الماء]\label{ex:g10:refraction:speed}
في الماء، $v = \dfrac{c}{n} = \dfrac{\qty{3.00e8}{m/s}}{1.33}
= \qty{2.26e8}{m/s}$. وفي الماس، يزحف الضوء بسرعة
$\qty{3.00e8}{m/s}/2.42 = \qty{1.24e8}{m/s}$ --- أي أقل من نصف سرعته
في الخلاء. أما الهواء فإنّ $n = 1.0003$ يغيّر السرعة بمقدار $0.03\%$: وفي
هذا الفصل نعامل الهواء معاملة الخلاء ونأخذ $n_{\text{الهواء}} = 1.00$.
\end{example}

\begin{theorem}[قانون سنيل للانكسار]\label{thm:g10:refraction:snell}
\index{قانون سنيل}
يقع \omterm{def:g10:refraction:refraction}{الشعاع المنكسر} في مستوي الشعاع الوارد \omterm{def:g10:refraction:angles}{والناظم}،
على الجهة الأخرى من \omterm{def:g10:refraction:angles}{الناظم}، وتحقق زاويتا الورود
\omterm{def:g10:refraction:refraction}{والانكسار} العلاقة
\[
n_1 \sin i_1 = n_2 \sin i_2 ,
\]
حيث $n_1$ و$n_2$ قرينتا \omterm{def:g10:refraction:refraction}{انكسار} الوسطين. والمسار
عكوس: فالضوء المنتقل رجوعًا على \omterm{def:g10:refraction:refraction}{الشعاع المنكسر}
يخرج على الشعاع الوارد.
\end{theorem}
\admitted

\begin{remark}\label{rem:g10:refraction:honest}
يُذكر قانونا الانعكاس \omterm{def:g10:refraction:refraction}{والانكسار} هنا بوصفهما واقعتين تجريبيتين.
ويُشتقّان اشتقاقًا أمينًا في مجلد السنة الأولى، حيث يُعامل الضوء
موجةً: فينتج القانونان --- والصيغة $n = c/v$ نفسها ---
عندئذٍ عن تباطؤ الموجة في الوسط الأكثف.
\end{remark}

\begin{omfigure}
\begin{tikzpicture}[scale=1.0, font=\small]
  \fill[blue!8] (-3.2,-2.1) rectangle (3.2,0);
  \draw[thick] (-3.2,0) -- (3.2,0);
  \draw[dashed] (0,-2.1) -- (0,2.3);
  \coordinate (O) at (0,0);
  \coordinate (I) at (-1.8,2.0);
  \coordinate (R) at (1.8,2.0);
  \coordinate (T) at (1.0,-1.73);
  \coordinate (N) at (0,2.3);
  \coordinate (Nd) at (0,-2.1);
  \draw[-Stealth, omDef, thick] (I) -- (O);
  \draw[-Stealth, omExo, thick] (O) -- (R);
  \draw[-Stealth, omThm, thick] (O) -- (T);
  \pic[draw, angle radius=9mm, "$i_1$", angle eccentricity=1.35] {angle = N--O--I};
  \pic[draw, angle radius=7mm, "$i_1'$", angle eccentricity=1.4] {angle = R--O--N};
  \pic[draw, angle radius=9mm, "$i_2$", angle eccentricity=1.35] {angle = Nd--O--T};
  \node[anchor=west] at (-3.1,0.45) {الوسط 1 ($n_1$)};
  \node[anchor=west] at (-3.1,-0.55) {الوسط 2 ($n_2$)};
\end{tikzpicture}

{\small الشعاع الوارد (بالأزرق) والمنعكس (بالرمادي) والمنكسر (بالأحمر). وتُقاس
كل الزوايا من \omterm{def:g10:refraction:angles}{الناظم} (المتقطع)؛ وعند دخول الوسط الأكثف
($n_2 > n_1$)، ينحني الشعاع \emph{نحو} \omterm{def:g10:refraction:angles}{الناظم}.}
\end{omfigure}

\begin{example}[من الهواء إلى الماء]\label{ex:g10:refraction:airwater}
يصيب شعاع بركة ساكنة بزاوية $i_1 = 30^\circ$. وبأخذ $n_1 = 1.00$ و
$n_2 = 1.33$:
\[
\sin i_2 = \frac{n_1 \sin i_1}{n_2} = \frac{\sin 30^\circ}{1.33}
= \frac{0.500}{1.33} = 0.376,
\qquad
i_2 = 22.1^\circ .
\]
فينحني الشعاع نحو \omterm{def:g10:refraction:angles}{الناظم}، كما يفعل دائمًا عند دخوله وسطًا
أبطأ ($n_2 > n_1$). وعند مغادرة الماء، يجري الحساب نفسه
عكسًا فينحني الشعاع \emph{مبتعدًا} عن \omterm{def:g10:refraction:angles}{الناظم}.
\end{example}

\begin{method}[حلّ مسألة انكسار]\label{met:g10:refraction:method}
\begin{enumerate}
  \item ارسم السطح \omterm{def:g10:refraction:angles}{والناظم} والشعاع؛ وعيّن $n_1$
        (وسط الشعاع الوارد) و$n_2$.
  \item تحقّق من أن الزوايا المعطاة مقيسة من \omterm{def:g10:refraction:angles}{الناظم}؛
        وحوّلها إن كانت مقيسة من السطح.
  \item اكتب $n_1 \sin i_1 = n_2 \sin i_2$ واعزل الجيب
        المجهول؛ ثم استخرج الزاوية بدالة الجيب العكسية في الآلة الحاسبة، كما في
        مجلد الرياضيات.
  \item تحقّق من السلامة: ينحني الشعاع نحو \omterm{def:g10:refraction:angles}{الناظم} إذا كان $n_2 > n_1$،
        ويبتعد عنه إذا كان $n_2 < n_1$، ويجب أن تخرج $\sin i_2$
        بين $0$ و$1$.
\end{enumerate}
\end{method}

والتجربة وراء القانون تسع طاولة مكتب: نصف قرص من البلكسيغلاس،
ومنقلة، وليزر. وتبدو أزواج الزوايا الخام $(i_1, i_2)$ غير منتظمة ---
وقد جدولها بطليموس نحو سنة 150 واستنتج خطأً أن
$i_2$ يتناسب مع $i_1$. أما إذا رسمت الجيب بدلالة الجيب،
استقامت المعطيات في خط يمرّ بالمبدأ ميله $n_1/n_2$: وذلك
الخط \emph{هو} \omterm{thm:g10:refraction:snell}{قانون سنيل}.

\begin{omfigure}
\begin{tikzpicture}
\begin{axis}[omaxis, width=7.2cm, height=5.4cm,
  xlabel={$\sin i_1$}, ylabel={$\sin i_2$},
  xmin=0, xmax=1.08, ymin=0, ymax=0.88,
  xtick={0,0.2,0.4,0.6,0.8,1.0}, ytick={0.2,0.4,0.6,0.8}]
  \addplot[omThm, thick, domain=0:1.02] {x/1.33};
  \addplot[only marks, mark=*, mark size=1.6pt, omDef] coordinates
    {(0.174,0.131) (0.342,0.257) (0.500,0.376) (0.643,0.483)
     (0.766,0.576) (0.866,0.651) (0.940,0.707)};
\end{axis}
\end{tikzpicture}

{\small قياسات من الهواء إلى الماء من $i_1 = 10^\circ$ إلى $70^\circ$:
إذا رُسم الجيب بدلالة الجيب اصطفّت المعطيات على خط يمرّ بالمبدأ
ميله $1/1.33$.}
\end{omfigure}

\begin{proposition}[العمق الظاهري]\label{prop:g10:refraction:depth}
\index{العمق الظاهري}
جسم على العمق $d$ تحت الماء، إذا نُظر إليه من فوق بزوايا صغيرة،
بدا على \emph{\omterm{prop:g10:refraction:depth}{العمق الظاهري}}
\[
d' = \frac{d}{n} ,
\]
حيث $n$ \omterm{def:g10:refraction:index}{قرينة انكسار} الماء. فمسبح عمقه \qty{1.20}{m} يبدو
عمقه $1.20/1.33 = \qty{0.90}{m}$ لا غير --- وهو سبب كلاسيكي
للغطسات الطائشة.
\end{proposition}
\admitted

\begin{remark}\label{rem:g10:refraction:depthhonest}
تنتج الصيغة عن تطبيق \omterm{thm:g10:refraction:snell}{قانون سنيل} على الأشعة شبه العمودية
الواصلة إلى العين؛ ويُجرى الاشتقاق في مجلد السنة الأولى. أما
اتجاه الأثر فلا يحتاج صيغة: فالأشعة المغادرة للماء تنحني
مبتعدة عن \omterm{def:g10:refraction:angles}{الناظم}، فتتعقبها العين رجوعًا إلى نقطة أعلى
من الجسم.
\end{remark}

\section{التبدد: قانون لكل لون}

يحوّل الموشور ضوء الشمس الأبيض إلى قوس قزح، كما بيّنت أطياف
\cref{ch:g10:light-spectra}. ويفسّر \omterm{def:g10:refraction:refraction}{الانكسار} \emph{لماذا}:
\omterm{def:g10:refraction:index}{فقرينة انكسار} وسطٍ ما ليست هي نفسها تمامًا لكل لون.

\begin{definition}[التبدد]\label{def:g10:refraction:dispersion}
يكون الوسط \emph{مبدِّدًا}\index{تبدد} إذا توقفت قرينة
انكساره على طول موجة الضوء. ففي الزجاج والماء تكون $n$
أكبر قليلًا للضوء الأزرق (طول موجة قصير) منها للضوء الأحمر
(طول موجة طويل) --- ولذلك ينحني الأزرق أكثر.
\end{definition}

\begin{example}[لونان وسطح واحد]\label{ex:g10:refraction:twocolors}
تصيب حزمة من \omterm{def:g10:light-spectra:white}{الضوء الأبيض} زجاج الكراون بزاوية $i_1 = 60^\circ$. ولهذا
الزجاج $n_{\text{الأحمر}} = 1.51$ و$n_{\text{الأزرق}} = 1.53$. \omterm{thm:g10:refraction:snell}{وقانون سنيل}،
مرة لكل لون:
\[
\sin i_{\text{الأحمر}} = \frac{\sin 60^\circ}{1.51} = 0.574,
\quad i_{\text{الأحمر}} = 35.0^\circ ;
\qquad
\sin i_{\text{الأزرق}} = \frac{\sin 60^\circ}{1.53} = 0.566,
\quad i_{\text{الأزرق}} = 34.5^\circ .
\]
فسطح واحد يفصل \omterm{def:g10:light-spectra:white}{الضوء الأبيض} بنصف درجة --- وهو غير مرئي على
لوح نافذة، لكن الموشور يكسر الضوء مرتين في جهة الدوران نفسها،
فيوسّع المروحة، وعلى بعد أمتار قليلة تفترق الألوان بسنتيمترات.
\end{example}

\begin{remark}\label{rem:g10:refraction:rainbow}
وقوس قزح هو الحساب نفسه مُجرًى داخل قطرة مطر: فضوء الشمس
ينكسر داخلًا القطرة، وينعكس على جدارها الخلفي، ثم ينكسر
مرة أخرى في طريق الخروج. وتمتد قرينة الماء من $1.331$ (أحمر) إلى $1.343$
(أزرق)، فيخرج كل لون بزاويته الخاصة --- ولا يرى راصدان
قوس قزح نفسه.
\end{remark}

\section{الانعكاس الكلي الداخلي}

في \omterm{thm:g10:refraction:snell}{قانون سنيل} فخّ. فبالانتقال من الماء نحو الهواء، ينحني \omterm{def:g10:refraction:refraction}{الشعاع
المنكسر} مبتعدًا عن \omterm{def:g10:refraction:angles}{الناظم}: $\sin i_2 = 1.33 \sin i_1 > \sin i_1$.
وبزيادة $i_1$، تطلب الصيغة سريعًا $\sin i_2 > 1$، وهو ما لا تعطيه
زاوية. فيكفّ \omterm{def:g10:refraction:refraction}{الانكسار} عندئذٍ عن الحدوث ببساطة.

\begin{definition}[الزاوية الحدية والانعكاس الكلي الداخلي]\label{def:g10:refraction:critical}
ليكن ضوء ينتقل في وسط قرينته $n_1$ نحو وسط قرينته
أصغر $n_2 < n_1$. و\emph{الزاوية الحدية}\index{الزاوية الحدية}
$i_c$ هي \omterm{def:g10:refraction:angles}{زاوية الورود} التي عندها يلامس \omterm{def:g10:refraction:refraction}{الشعاع المنكسر}
السطح ($i_2 = 90^\circ$). وإذا كان $i_1 > i_c$ فلا وجود \omterm{def:g10:refraction:refraction}{لشعاع منكسر}:
إذ يعكس السطح كل الضوء راجعًا إلى الوسط الأول، وفق
قانون الانعكاس. وهذا هو \emph{الانعكاس الكلي
الداخلي}\index{الانعكاس الكلي الداخلي}.
\end{definition}

\begin{proposition}[صيغة الزاوية الحدية]\label{prop:g10:refraction:criticalformula}
إذا كان $n_1 > n_2$، فإن
\[
\sin i_c = \frac{n_2}{n_1} .
\]
\end{proposition}

\begin{proof}
ضَع $i_2 = 90^\circ$ في \omterm{thm:g10:refraction:snell}{قانون سنيل}: $n_1 \sin i_c = n_2 \sin 90^\circ
= n_2$. وإذا كان $i_1 > i_c$، فسيقتضي \omterm{thm:g10:refraction:snell}{قانون سنيل}
$\sin i_2 = \frac{n_1}{n_2}\sin i_1 > 1$: أي لا وجود لاتجاه منكسر.
(أما كون الضوء ينعكس عندئذٍ \emph{كليًّا} --- لا يُمتصّ --- فواقعة
تجريبية، تؤكدها النظرية الموجية في مجلد السنة
الأولى.)
\end{proof}

\begin{omfigure}
\begin{tikzpicture}[scale=0.92, font=\small]
  \fill[blue!8] (-3.6,-1.9) rectangle (3.6,0);
  \draw[thick] (-3.6,0) -- (3.6,0);
  \draw[dashed] (-2.2,-1.6) -- (-2.2,1.6);
  \draw[dashed] (0.3,-1.6) -- (0.3,1.6);
  \draw[dashed] (2.2,-1.6) -- (2.2,1.6);
  \draw[-Stealth, omDef, thick] (-2.95,-1.30) -- (-2.2,0);
  \draw[-Stealth, omThm, thick] (-2.2,0) -- (-1.20,1.12);
  \draw[-Stealth, omDef, thick] (-0.83,-0.99) -- (0.3,0);
  \draw[-Stealth, omThm, thick] (0.3,0) -- (1.75,0.06);
  \draw[-Stealth, omDef, thick] (0.90,-0.75) -- (2.2,0);
  \draw[-Stealth, omDef, thick] (2.2,0) -- (3.50,-0.75);
  \node at (-2.2,-1.85) {$i_1 < i_c$};
  \node at (0.3,-1.85) {$i_1 = i_c$};
  \node at (2.2,-1.85) {$i_1 > i_c$};
\end{tikzpicture}

{\small من الماء نحو الهواء، بورود متزايد: \omterm{def:g10:refraction:refraction}{انكسار} مبتعد
عن \omterm{def:g10:refraction:angles}{الناظم}، ثم \omterm{def:g10:refraction:refraction}{شعاع منكسر} يلامس السطح عند
\omterm{def:g10:refraction:critical}{الزاوية الحدية}، ثم انعكاس كلي داخلي --- فقد صار السطح
مرآة تامة.}
\end{omfigure}

\begin{example}[ثلاث زوايا حدية]\label{ex:g10:refraction:criticalvalues}
نحو الهواء: للماء $\sin i_c = 1/1.33 = 0.752$، ومنه
$i_c = 48.8^\circ$؛ وللزجاج $\sin i_c = 1/1.50$، ومنه
$i_c = 41.8^\circ$؛ وللماس $\sin i_c = 1/2.42 = 0.413$، ومنه
$i_c = 24.4^\circ$. فكلما ارتفعت القرينة ضاق مخروط الإفلات:
ففي الماس، كل شعاع يبعد أكثر من $24.4^\circ$ عن \omterm{def:g10:refraction:angles}{ناظم} وجه من الأوجه
يبقى محبوسًا --- وذلك سرّ بريقه، وهو مشرَّح في
\cref{pb:g10:refraction:1}.
\end{example}

\section{الألياف البصرية}

يحوّل \omterm{def:g10:refraction:critical}{الانعكاس الكلي الداخلي} خيطًا زجاجيًّا إلى أنبوب للضوء:
فالشعاع الداخل شبه الموازي للمحور يصيب الجدار بزاوية أبعد كثيرًا
من \omterm{def:g10:refraction:critical}{الزاوية الحدية}، فينعكس دون ضياع، ويفعل ذلك ملايين المرات
في الكيلومتر دون تسرّب.

\begin{definition}[الليف البصري]\label{def:g10:refraction:fiber}
\emph{الليف البصري}\index{ليف بصري} خيط زجاجي رفيع مصنوع
من \emph{قلب}\index{قلب} قرينته $n_1$ يلفّه
\emph{غلاف}\index{غلاف} قرينته أصغر قليلًا $n_2 < n_1$.
والضوء المحقون في أحد طرفيه ينتقل على طول القلب بانعكاس كلي
داخلي متكرر على الحد الفاصل بين القلب والغلاف.
\end{definition}

\begin{omfigure}
\begin{tikzpicture}[scale=1.0, font=\small]
  \draw[omExo, thick] (0,0.85) -- (7,0.85);
  \draw[omExo, thick] (0,-0.85) -- (7,-0.85);
  \fill[blue!8] (0,-0.5) rectangle (7,0.5);
  \draw[thick] (0,0.5) -- (7,0.5);
  \draw[thick] (0,-0.5) -- (7,-0.5);
  \draw[-Stealth, omThm, thick] (0,-0.36) -- (1.0,0.5) -- (2.16,-0.5)
    -- (3.32,0.5) -- (4.48,-0.5) -- (5.64,0.5) -- (6.6,-0.33);
  \node[anchor=south] at (3.5,0.9) {الغلاف ($n_2 = 1.46$)};
  \node[fill=blue!8, inner sep=2pt] at (3.5,-0.25) {القلب ($n_1 = 1.48$)};
\end{tikzpicture}

{\small ضوء موجَّه على طول \omterm{def:g10:refraction:fiber}{قلب} ليف بانعكاس كلي
داخلي (وزوايا الارتداد مبالَغ فيها كثيرًا: فالأشعة الحقيقية تلامس الجدار
بأكثر من $80^\circ$ عن \omterm{def:g10:refraction:angles}{الناظم}).}
\end{omfigure}

\begin{example}[ليف بالأرقام]\label{ex:g10:refraction:fibernumbers}
بأخذ $n_1 = 1.48$ و$n_2 = 1.46$:
\[
\sin i_c = \frac{1.46}{1.48} = 0.986,
\qquad i_c = 80.6^\circ :
\]
فلا تُوجَّه إلا الأشعة الواقعة في حدود $9^\circ$ عن المحور --- وهذا
مناسب، إذ إن ذلك بالضبط هو أسلوب حقن ضوء الليزر. وداخل
\omterm{def:g10:refraction:fiber}{القلب} ينتقل الضوء بسرعة $v = c/1.48 = \qty{2.03e8}{m/s}$: فتعبر نبضة
ليفًا طوله \qty{100}{km} في نحو نصف مِلّي \omterm{def:g10:orders-of-magnitude:unit}{ثانية}.
\end{example}

\begin{remark}[رتبة قدر وصلات المعطيات]\label{rem:g10:refraction:datarates}
تحمل الكابلات البحرية --- وهي حزم من بضع عشرات الألياف على قاع المحيط ---
كل الحركة العابرة للقارات تقريبًا. فالليف الحديث الواحد
ينقل من رتبة $10^{13}$ بت في \omterm{def:g10:orders-of-magnitude:unit}{الثانية} بإرسال أطوال موجة
كثيرة في آن؛ ويبلغ الكابل الكامل $10^{14}$ بت في \omterm{def:g10:orders-of-magnitude:unit}{الثانية}،
مع مضخّم كل \qty{80}{km} أو نحو ذلك. أما الأقمار الاصطناعية، للمقارنة،
فتتولى جزءًا صغيرًا من واحد بالمئة من الحركة.
\end{remark}

\section{تمارين}

\begin{exercise}[$\star$]\label{exo:g10:refraction:1}
تصيب حزمة ليزر مرآة مستوية بزاوية $25^\circ$ \emph{من سطح
المرآة}. أعطِ \omterm{def:g10:refraction:angles}{زاوية الورود}، وزاوية الانعكاس، والزاوية
بين الحزمتين الواردة والمنعكسة.
\end{exercise}

\begin{exercise}[$\star$]\label{exo:g10:refraction:2}
\begin{enumerate}
  \item احسب سرعة الضوء في الماء ($n = 1.33$) وفي الماس
        ($n = 2.42$).
  \item في جسم صلب ما، ينتقل الضوء بسرعة \qty{1.97e8}{m/s}. احسب
        قرينة انكساره وعيّن المادة.
\end{enumerate}
\end{exercise}

\begin{exercise}[$\star$]\label{exo:g10:refraction:3}
يصيب شعاع في الهواء كتلة زجاجية ($n = 1.50$) بزاوية $i_1 = 40^\circ$.
احسب زاوية \omterm{def:g10:refraction:refraction}{الانكسار}. وبأي زاوية يجب أن يصيب شعاع منتقل
داخل الزجاج السطحَ ليخرج إلى الهواء بزاوية
$40^\circ$؟
\end{exercise}

\begin{exercise}[$\star$]\label{exo:g10:refraction:4}
يصل ضوء الشمس إلى سطح بحيرة بزاوية $55^\circ$ عن العمودي.
احسب زاوية \omterm{def:g10:refraction:refraction}{الشعاع المنكسر}، وقل أينحني الشعاع
نحو \omterm{def:g10:refraction:angles}{الناظم} أم مبتعدًا عنه --- مع سبب من كلمة واحدة.
\end{exercise}

\begin{exercise}[$\star$]\label{exo:g10:refraction:5}
لتعيين سائل صافٍ، يُرسل شعاع على سطحه بزاوية
$i_1 = 45^\circ$؛ فيُقاس \omterm{def:g10:refraction:refraction}{الشعاع المنكسر} عند $i_2 = 32^\circ$.
احسب \omterm{def:g10:refraction:index}{قرينة الانكسار} وعيّن السائل من جدول
\cref{def:g10:refraction:index}.
\end{exercise}

\begin{exercise}[$\star\star$]\label{exo:g10:refraction:6}
يعطي نصف قرص من بلاستيك شفاف القياسات التالية:
\begin{center}
\begin{tabular}{lccccc}
$i_1$ & $20^\circ$ & $30^\circ$ & $40^\circ$ & $50^\circ$ & $60^\circ$ \\
$i_2$ & $13.5^\circ$ & $19.5^\circ$ & $25.5^\circ$ & $31.0^\circ$
      & $35.5^\circ$ \\
\end{tabular}
\end{center}
\begin{enumerate}
  \item احسب $\dfrac{\sin i_1}{\sin i_2}$ لكل زوج. فماذا
        تلاحظ؟
  \item استنتج \omterm{def:g10:refraction:index}{قرينة انكسار} البلاستيك وعيّنه.
  \item لماذا يكون رسم $\sin i_2$ بدلالة $\sin i_1$ اختبارًا \omterm{thm:g10:refraction:snell}{لقانون سنيل}
        أفضل من رسم $i_2$ بدلالة $i_1$؟
\end{enumerate}
\end{exercise}

\begin{exercise}[$\star\star$]\label{exo:g10:refraction:7}
احسب \omterm{def:g10:refraction:critical}{الزاوية الحدية} نحو الهواء للماء وللزجاج ($n = 1.50$) وللماس.
ثم اشرح، \omterm{thm:g10:refraction:snell}{بقانون سنيل}، لماذا لا توجد زاوية حدية
للضوء الذاهب \emph{من} الهواء \emph{إلى} الماء.
\end{exercise}

\begin{exercise}[$\star\star$]\label{exo:g10:refraction:8}
\begin{enumerate}
  \item عمق مسبح \qty{2.0}{m}. فأي عمق يدركه سابح واقف
        على الحافة؟
  \item يرقب مالك الحزين سمكة تسبح على عمق \qty{40}{cm} تحت السطح. فعلى
        أي عمق تبدو السمكة، وهل ينبغي أن ينقضّ الطائر
        فوق الصورة التي يراها أم عليها أم تحتها؟
\end{enumerate}
\end{exercise}

\begin{exercise}[$\star\star$]\label{exo:g10:refraction:9}
يصيب شعاع لوح نافذة ($n = 1.50$، وجهاه متوازيان) بزاوية $50^\circ$.
\begin{enumerate}
  \item احسب زاوية \omterm{def:g10:refraction:refraction}{الانكسار} عند الوجه الأول.
  \item يصيب \omterm{def:g10:refraction:refraction}{الشعاع المنكسر} الوجه الثاني من الداخل، بالزاوية
        نفسها. احسب زاوية الخروج وقارنها بالزاوية $50^\circ$.
  \item فماذا تفعل إذن نافذة سميكة باتجاه شعاع يعبرها
        وبموضعه؟
\end{enumerate}
\end{exercise}

\begin{exercise}[$\star\star$]\label{exo:g10:refraction:10}
يصيب \omterm{def:g10:light-spectra:white}{ضوء أبيض} كتلة من زجاج الفلينت بزاوية $55^\circ$. ولهذا
الزجاج $n_{\text{الأحمر}} = 1.61$ و$n_{\text{الأزرق}} = 1.66$. احسب
زاويتي \omterm{def:g10:refraction:refraction}{انكسار} الشعاعين الأحمر والأزرق والزاوية بينهما.
وأيّ اللونين ينتهي أقرب إلى \omterm{def:g10:refraction:angles}{الناظم}؟
\end{exercise}

\begin{exercise}[$\star\star$]\label{exo:g10:refraction:11}
\omterm{def:g10:refraction:fiber}{لليف بصري} \omterm{def:g10:refraction:fiber}{قلب} قرينته $1.48$ \omterm{def:g10:refraction:fiber}{وغلاف} قرينته
$1.46$.
\begin{enumerate}
  \item احسب \omterm{def:g10:refraction:critical}{الزاوية الحدية} على الحد الفاصل بين \omterm{def:g10:refraction:fiber}{القلب} \omterm{def:g10:refraction:fiber}{والغلاف}.
  \item يرتدّ شعاع موجَّه عند \omterm{def:g10:refraction:critical}{الزاوية الحدية} بالضبط. بيّن
        أنه يقطع على مدى \qty{1.0}{km} من الليف المستقيم نحو
        \qty{14}{m} زيادةً على ما يقطعه شعاع ذاهب على المحور مباشرة.
\end{enumerate}
\end{exercise}

\begin{exercise}[$\star\star\star$]\label{exo:g10:refraction:12}
عين غطّاس على عمق \qty{2.0}{m} تحت سطح ساكن تمامًا. وبسبب
\omterm{def:g10:refraction:refraction}{الانكسار}، يرى الغطّاس السماء \emph{كلها} مضغوطة في
قرص مضيء فوق رأسه مباشرة --- هي نافذة سنيل.
\begin{enumerate}
  \item اشرح لماذا يصل الضوء الآتي من الأفق إلى عين الغطّاس
        \omterm{def:g10:refraction:critical}{بالزاوية الحدية} للماء، أي $48.8^\circ$ عن العمودي.
  \item احسب نصف قطر القرص المضيء على السطح، مستعملًا
        الظل في المثلث القائم: العين، العمودي، حافة القرص.
  \item فماذا يرى الغطّاس على السطح \emph{خارج} القرص؟
\end{enumerate}
\end{exercise}

\begin{exercise}[$\star\star\star$]\label{exo:g10:refraction:13}
لموشور زوايا $30^\circ$ و$60^\circ$ و$90^\circ$. تدخل حزمة رفيعة من
\omterm{def:g10:light-spectra:white}{الضوء الأبيض} عمودية على أحد الضلعين المجاورين للزاوية القائمة،
فتعبر دون انحراف، وتصيب الوجه الطويل من الداخل بزاوية
$i_1 = 30^\circ$. ولهذا الزجاج $n_{\text{الأحمر}} = 1.51$ و
$n_{\text{الأزرق}} = 1.53$.
\begin{enumerate}
  \item احسب زاويتي خروج الشعاعين الأحمر والأزرق والعرض
        الزاوي للمروحة بينهما.
  \item يُستبدل بالموشور موشور $45^\circ$--$45^\circ$--$90^\circ$،
        فتصير \omterm{def:g10:refraction:angles}{زاوية الورود} الداخلية $45^\circ$.
        بيّن أنه لا يخرج أي ضوء عبر الوجه الطويل البتة.
\end{enumerate}
\end{exercise}

\begin{exercise}[$\star\star\star$]\label{exo:g10:refraction:14}
تطوي المِرقابات الضوء بمواشير زجاجية $45^\circ$--$45^\circ$--$90^\circ$
($n = 1.50$): إذ تدخل الحزمة أحد الوجهين القصيرين عموديًّا
وتصيب الوجه الطويل من الداخل بزاوية $45^\circ$.
\begin{enumerate}
  \item بيّن أن الوجه الطويل يعمل عمل مرآة تامة. ولماذا يكون هذا
        أفضل من مرآة معدنية تمتصّ بضعة بالمئة عند
        كل انعكاس؟
  \item يغرق المِرقاب فيصير الوجه الطويل على تماس مع
        الماء. احسب \omterm{def:g10:refraction:critical}{الزاوية الحدية} الجديدة للحد الفاصل بين الزجاج والماء
        وبيّن أن المرآة تخفق.
  \item بأي زاوية يغادر ضوءُ السؤال 2 الموشورَ إلى
        الماء؟
\end{enumerate}
\end{exercise}

\begin{exercise}[$\star\star\star$]\label{exo:g10:refraction:15}
يمتد كابل عابر للأطلسي على \qty{6200}{km} من الليف (قرينة \omterm{def:g10:refraction:fiber}{القلب} $1.47$)
بين لندن ونيويورك.
\begin{enumerate}
  \item احسب سرعة الضوء في \omterm{def:g10:refraction:fiber}{القلب} وزمن انتقال نبضة في
        اتجاه واحد.
  \item تنتقل موجة راديوية بسرعة $c$. فكم تحتاج على
        المسافة نفسها، وكم يكلّف زجاج الليف من زمن الذهاب والإياب مقارنةً بالراديو؟
  \item تدفع بعض الشركات المالية ثروات لاقتطاع أجزاء من الألف من \omterm{def:g10:orders-of-magnitude:unit}{الثانية} من هذه
        الوصلة. اقترح طريقتين فيزيائيتين متمايزتين لجعل الضوء
        يصل أبكر.
\end{enumerate}
\end{exercise}

\section{مسألة: الماسة والمسبح والليف}

\begin{problem}[{مسألة نهاية الأسبوع --- ثلاث حيوات لقانون سنيل: مسبح
يكذب على عمقه، وحجر مقطوع بحيث لا يقدر الضوء على مغادرته،
والمحيط ذو الستين جزءًا من الألف من الثانية تحت الإنترنت}]
\label{pb:g10:refraction:1}
مساواة واحدة بين جيبين، وثلاث مِهَن. ففي مسبح تخدع
عينيك في العمق وتعصر السماء كلها في دائرة. وفي
ماسة، مدفوعةً إلى حدّها الأقصى، ترفض أن تدع الضوء يخرج ---
وقد قطع الجواهريون الأحجار حول ذلك الرفض قرونًا. وتحت
الأطلسي توجّه نبضات الليزر داخل خيط من الزجاج، وتضع
حدّ سرعة الإنترنت العالمي. وتعالج هذه المسألة
الحيوات الثلاث بالترتيب، بالقانون نفسه في كل مرة.

\medskip
\noindent\textbf{الجزء I --- الإحماء.}
\begin{enumerate}
  \item احسب سرعة الضوء في الماء ($n = 1.33$) وفي الماس
        ($n = 2.42$). بأي عامل يبطّئ الماس الضوء؟
  \item يصيب شعاع في الهواء بركة بزاوية $i_1 = 40^\circ$. احسب
        زاوية \omterm{def:g10:refraction:refraction}{الانكسار}.
  \item يصيب شعاع في الماء السطح من تحت بزاوية $28.9^\circ$.
        فبأي زاوية يخرج إلى الهواء؟ واذكر المبدأ
        الذي يوضّحه ذلك.
  \item احسب \omterm{def:g10:refraction:critical}{الزاوية الحدية} للحد الفاصل بين الماء والهواء، ثم
        صِف ما يحدث لشعاع تحت الماء يصيب السطح
        بزاوية $60^\circ$.
  \item قرينة الهواء $1.0003$. احسب سرعة الضوء في الهواء
        وبرّر التقريب $n_{\text{الهواء}} = 1.00$ المستعمل
        في كل مكان.
\end{enumerate}

\medskip
\noindent\textbf{الجزء II --- المسبح.}
\begin{enumerate}[resume]
  \item عمق مسبح \qty{3.0}{m}. باستعمال صيغة \omterm{prop:g10:refraction:depth}{العمق الظاهري}
        (\cref{prop:g10:refraction:depth})، احسب العمق الذي يدركه شخص
        واقف فوقه.
  \item اشرح برسم بشعاعين (موصوف بالكلمات) لماذا تبدو قطعة نقود
        في القاع مرتفعة: ماذا تفعل الأشعة المغادرة للماء،
        وأين تضع العين تقاطعها؟
  \item عين غطّاس على \qty{3.0}{m}. احسب نصف قطر نافذة
        سنيل --- أي القرص المضيء الذي يرى الغطّاس من خلاله
        السماء كلها (\cref{exo:g10:refraction:12}).
  \item يرسل مصباح في القاع أشعة إلى أعلى بزوايا $30^\circ$
        و$45^\circ$ و$50^\circ$ عن العمودي. لكل شعاع، قل
        أيفلت وبأي زاوية، أم ماذا يحدث له بدل ذلك.
  \item يصوّب صياد بالحربة على الضفة نحو سمكة. أينبغي أن تكون الضربة
        فوق صورة السمكة أم عليها أم تحتها --- وهل ترى
        السمكة الصيادَ مزاحًا هو أيضًا حين تنظر إليه؟
\end{enumerate}

\medskip
\noindent\textbf{الجزء III --- الماسة.}
\begin{enumerate}[resume]
  \item احسب \omterm{def:g10:refraction:critical}{الزاوية الحدية} للماس ($n = 2.42$) نحو الهواء،
        وقارنها بزاوية الزجاج ($n = 1.52$)،
        وهي $41.1^\circ$.
  \item يصيب شعاع داخل الحجر وجهًا بزاوية $30^\circ$ عن
        ناظمه. فماذا يحدث في الماس؟ وفي التقليد الزجاجي؟
        اشرح في جملة واحدة لماذا ينتج عن قطع الماسة، مقرونًا
        بزاويتها الحدية الصغيرة، ذلك البريق.
  \item والماس مبدِّد بقوة أيضًا: $n_{\text{الأحمر}} = 2.41$،
        $n_{\text{الأزرق}} = 2.45$. يدخل \omterm{def:g10:light-spectra:white}{ضوء أبيض} وجهًا بزاوية
        $45^\circ$؛ احسب زاوية \omterm{def:g10:refraction:refraction}{الانكسار} لكل لون والفصل
        الزاوي.
  \item ثم يبلغ الشعاعان الأحمر والأزرق كلاهما وجه خروج من الداخل
        بزاوية $17.0^\circ$. احسب زاويتي الخروج وفصل
        المروحة الخارجة. وقارنه بالفصل عند الدخول: فماذا يحدث
        \omterm{def:g10:refraction:dispersion}{للتبدد} قرب \omterm{def:g10:refraction:critical}{الزاوية الحدية}؟
  \item تغطي قطرة ماء ($n = 1.33$) أحد الأوجه. احسب \omterm{def:g10:refraction:critical}{الزاوية
        الحدية} الجديدة للحد الفاصل بين الماس والماء، وحدّد مصير
        الشعاع $30^\circ$ في السؤال 12، واشرح لماذا يبقي الجواهريون
        الماس نظيفًا بدقة.
\end{enumerate}

\medskip
\noindent\textbf{الجزء IV --- الليف.}
\begin{enumerate}[resume]
  \item لليف \omterm{def:g10:refraction:fiber}{قلب} قرينته $1.48$ \omterm{def:g10:refraction:fiber}{وغلاف} قرينته
        $1.46$. احسب \omterm{def:g10:refraction:critical}{الزاوية الحدية} على الحد الفاصل بين \omterm{def:g10:refraction:fiber}{القلب}
        \omterm{def:g10:refraction:fiber}{والغلاف}.
  \item خيط زجاجي عارٍ في الهواء تكون زاويته الحدية
        $42.5^\circ$ --- وهو أفضل بكثير في الظاهر. أعطِ سببين لاستعمال
        \omterm{def:g10:refraction:fiber}{الغلاف} مع ذلك. (فكّر في الغبار والخدوش، وفيما يفعله
        لمس السطح \omterm{def:g10:refraction:critical}{بالانعكاس الكلي الداخلي}.)
  \item احسب سرعة الضوء في \omterm{def:g10:refraction:fiber}{القلب}، ثم زمن انتقال
        نبضة في اتجاه واحد على \qty{6000}{km} من الليف.
  \item شعاع يرتدّ عند \omterm{def:g10:refraction:critical}{الزاوية الحدية} بالضبط يقطع
        $1/\sin i_c$ ضعف طول المحور. احسب المسافة
        الزائدة على \qty{6000}{km} والزمن الزائد الذي تكلّفه.
  \item بيت القصيد. يحمل زوج ليف حديث نحو $10^{13}$ بت
        في \omterm{def:g10:orders-of-magnitude:unit}{الثانية}، وهذا الكتاب نحو $2.4 \times 10^8$ بت.
        احسب زمن الذهاب والإياب (``البِنغ'') عبر المحيط وعدد
        نسخ هذا الكتاب التي ينقلها الليف في \omterm{def:g10:orders-of-magnitude:unit}{الثانية} ---
        ثم قل، في جملة واحدة، ماذا يفعل \omterm{thm:g10:refraction:snell}{قانون سنيل}
        للإنترنت.
\end{enumerate}
\end{problem}
